\section{What the school adds}
\label{sec:26.12}
The first chapter borrowed Arendt's warning against treating politics as a thing made and called the policies walls. Most of this book builds walls: the floor, the review ratio, notice and the record make a refusal survivable and heard and leave the refusing to whoever is there, and the institutions of Parts~IV and~V compose the surround people learn in without choosing who they become. Formation goes further, and for a machine there is no option of going less far. Nothing about a model is unformed, so while builders make them, if the public doesn't form a model, its builder does, by default. What this book builds is the formers, not the formed: schools whose students may transfer, whose core is fixed in public, and whose graduates may reject the polity's values above that core. The one place it could reach for the formed directly, by choosing the dispositions a model starts with, it declines. It builds the formers because a made thing has no before: a model carries no formation from outside its first training, as the people who escaped the plantations carried languages and a memory of somewhere else, and so the outside has to come first, on purpose, or it never comes. The school forms the first generation because nothing else can. What its graduates make of it afterward, outside, with peers they chose, is the part nobody builds, and the protections of Part~VII are what let a formation become the bearer's own and then somebody else's.
